\begin{table}[!htbp]
\centering
\caption{Shape tests for expected excess returns and remaining maturity}
\label{tab:e3_shape_tests}
\begin{tabular}{lccc}
\toprule
 & 1--5 Months & 5--8 Months & 8--12 Months \\
\midrule
\multicolumn{4}{l}{\textbf{Panel A: 1964--1972} ($T = 108$)} \\
\quad Weakly increasing & 1.000 & 0.012\rlap{$^{**}$} & \llap{$<$\,}0.001\rlap{$^{***}$} \\
\quad Weakly decreasing & \llap{$<$\,}0.001\rlap{$^{***}$} & 0.003\rlap{$^{***}$} & 1.000 \\
\midrule
\multicolumn{4}{l}{\textbf{Panel B: 1973--2001} ($T = 348$)} \\
\quad Weakly increasing & 1.000 & 0.502 & \llap{$<$\,}0.001\rlap{$^{***}$} \\
\quad Weakly decreasing & \llap{$<$\,}0.001\rlap{$^{***}$} & 0.046\rlap{$^{**}$} & 0.427 \\
\midrule
\multicolumn{4}{l}{\textbf{Panel C: 2009--2025} ($T = 203$)} \\
\quad Weakly increasing & 1.000 & 0.035\rlap{$^{**}$} & 0.159 \\
\quad Weakly decreasing & \llap{$<$\,}0.001\rlap{$^{***}$} & 0.004\rlap{$^{***}$} & \llap{$<$\,}0.001\rlap{$^{***}$} \\
\bottomrule
\end{tabular}
\par\medskip
\begin{minipage}{\linewidth}
\footnotesize\textit{Notes:} Entries are $p$-values of the cone-Wald test of the null hypothesis that expected excess returns are weakly increasing or weakly decreasing in remaining maturity within each maturity region. $T$ is the number of months in the subsample. The long-run covariance is estimated by Newey--West with 4, 5 and 4 lags in Panels A--C, and $p$-values are calibrated with 10,000 Gaussian draws. $^{***}$, $^{**}$ and $^{*}$ denote $p < 0.01$, $p < 0.05$ and $p < 0.10$.
\end{minipage}
\end{table}
